\begin{proof}[Proof of Theorem~\ref{thm:14.6i} from Theorem~\ref{thm:14.6ii}]
	We proceed by induction on $k$, observing that the $k=0$ case of Theorem~\ref{thm:14.6i} follows immediately from the $k=0$ case of Theorem~\ref{thm:14.6ii}. 
	Now assume $k>0$, and fix $\mu\vdash mn$ with $l(\mu)=n-k$. Let $\nu=\mu\sqcup(1^k)$, so $l(\nu)=n$ and $\hat{\nu}=\hat{\mu}$.
	
	We aim to evaluate $a^{\nu/(1^k)}_{\lambda',(m)}$, and compare it to Theorem~\ref{thm:14.6ii}. To do this, we will study the constituents in the skew character $\chi^{\nu/(1^k)}$. First, note that for any $\omega\vdash|\nu|-k$, by Theorem~\ref{thm:LR} we must have $c^\nu_{\omega,(1^k)}\in\{0,1\}$. Moreover, $c^\nu_{\omega,(1^k)}=1$ if and only if $\omega\subseteq\nu$ and all $k$ boxes of $[\nu/\omega]$ belong to different rows. We will denote by $\mathcal{A}$ the collection of $\omega$ with $c^\nu_{\omega,(1^k)}=1$. In particular,
	\[ \sum_{\omega\in\mathcal{A}}\chi^\omega = \chi^{{\nu}/(1^k)}. \]
	We partition $\mathcal{A}$ as a disjoint union, $\mathcal{A} = \coprod_{j=0}^{k}\mathcal{A}_j$, where $\mathcal{A}_j$ is the collection of $\omega\in\mathcal{A}$ for which $l(\omega)=n-k+j$. Notice that for each $j$, $\mathcal{A}_j$ bijects to $\mathcal{B}_j \coloneqq  \{\varpi\vdash|\hat{\mu}|-j \mid c^{\hat{\mu}}_{\varpi,(1^j)}\}$; the map is given by removal of the first column, $\omega\mapsto\hat{\omega}$, and this is seen to be a bijection by Theorem~\ref{thm:LR}. By a similar application of Theorem~\ref{thm:LR},
	
	\begin{equation}\label{eqn:B_j}
		\sum_{\varpi\in\mathcal{B}_j}\chi^\varpi = \chi^{\hat{\mu}/(1^j)}.
	\end{equation}
	Now observe that
	\[ a^{\nu/(1^k)}_{\lambda',(m)} = \sum_{\omega\in\mathcal{A}} a^{\omega}_{\lambda',(m)} = a^{\mu}_{\lambda',(m)} + \sum_{j=1}^k \sum_{\omega\in\mathcal{A}_j} a^{\omega}_{\lambda',(m)}. \]
	The idea here is that in our partition of $\mathcal{A}$, the $j=0$ term contributes precisely $a^{\mu}_{\lambda',(m)}$. Let us set $X \coloneqq  -\sum_{j=1}^k \sum_{\omega\in\mathcal{A}_j} a^{\omega}_{\lambda',(m)}$. Assuming Theorem~\ref{thm:14.6ii}, it suffices to show that $X$ equals the $\sum_{i=0}^{k-1}(\dotsc)$ part of the summation on the right hand side of \eqref{eqn:A}. For $0<j\leq k$, we have that
	\begin{align*}
		\sum_{\omega\in\mathcal{A}_j} a^{\omega}_{\lambda',(m)} &\overset{\text{\normalfont\tiny(inductive hypothesis)}}{=} \sum_{\omega\in\mathcal{A}_j} \sum_{i=0}^{k-j}(-1)^{k-j+i}\sum_{\substack{\gamma\vdash k-j + (m-1)i\\ \beta\vdash i}}a^{\gamma/(k-j-i)}_{\beta',(m)}\cdot a^{\hat{\omega}/\gamma}_{\lambda/\beta,(m-1)}\\
		&\overset{\text{\normalfont\tiny(bijection $\mathcal{A}_j\to\mathcal{B}_j$)}}{=}\sum_{\varpi\in\mathcal{B}_j} \sum_{i=0}^{k-j}(-1)^{k-j+i}\sum_{\substack{\gamma\vdash k-j + (m-1)i\\ \beta\vdash i}}a^{\gamma/(k-j-i)}_{\beta',(m)}\cdot a^{\varpi/\gamma}_{\lambda/\beta,(m-1)}
		\\
		&\overset{\normalfont\tiny \text{\eqref{eqn:a}}}{=}\sum_{\varpi\in\mathcal{B}_j} \sum_{i=0}^{k-j}(-1)^{k-j+i}\sum_{\substack{\gamma\vdash k-j + (m-1)i\\ \beta\vdash i}}a^{\gamma/(k-j-i)}_{\beta',(m)}\cdot \left\langle\chi^{\varpi},\rho^{\lambda/\beta}_{m-1}\boxtimes\chi^\gamma\right\rangle \\
		&\overset{\normalfont\tiny\text{\eqref{eqn:B_j}}}{=}\sum_{i=0}^{k-j}(-1)^{k-j+i}\sum_{\substack{\gamma\vdash k-j + (m-1)i\\ \beta\vdash i}}a^{\gamma/(k-j-i)}_{\beta',(m)}\cdot \left\langle\chi^{\hat{\mu}/(1^j)},\rho^{\lambda/\beta}_{m-1}\boxtimes\chi^\gamma\right\rangle
		\\
		&=\sum_{i=0}^{k-j}(-1)^{k-j+i}\sum_{\substack{\gamma\vdash k-j + (m-1)i\\ \beta\vdash i}} a^{\gamma/(k-j-i)}_{\beta',(m)}\cdot \left\langle\chi^{\hat{\mu}},\rho^{\lambda/\beta}_{m-1}\boxtimes\chi^\gamma\boxtimes\chi^{(1^j)}\right\rangle\\
		&=\sum_{i=0}^{k-j}(-1)^{k-j+i}\sum_{\substack{\gamma\vdash k-j + (m-1)i\\ \alpha\vdash k+(m-1)i\\ \beta\vdash i}} c^{\alpha}_{\gamma,(1^j)}\cdot a^{\gamma/(k-j-i)}_{\beta',(m)}\cdot \left\langle\chi^{\hat{\mu}},\rho^{\lambda/\beta}_{m-1}\boxtimes\chi^\alpha\right\rangle\\
		&\overset{\normalfont\tiny\text{\eqref{eqn:a}}}{=}\sum_{i=0}^{k-j}(-1)^{k-j+i}\sum_{\substack{\gamma\vdash k-j + (m-1)i\\ \alpha\vdash k+(m-1)i\\ \beta\vdash i}} c^{\alpha}_{\gamma,(1^j)}\cdot a^{\gamma/(k-j-i)}_{\beta',(m)}\cdot a^{\hat{\mu}/\alpha}_{\lambda/\beta,(m-1)}.
	\end{align*}
	It follows that
	\begin{align*}
		X &= -\sum_{j=1}^k \sum_{i=0}^{k-j}(-1)^{k-j+i}\sum_{\substack{\gamma\vdash k-j + (m-1)i\\ \alpha\vdash k+(m-1)i\\ \beta\vdash i}} c^{\alpha}_{\gamma,(1^j)}\cdot a^{\gamma/(k-j-i)}_{\beta',(m)}\cdot a^{\hat{\mu}/\alpha}_{\lambda/\beta,(m-1)}\\
		&= \sum_{i=0}^{k-1} \sum_{j=1}^{k-i} (-1)^{k+i} \cdot (-1)^{j+1} \sum_{\substack{\gamma\vdash k-j + (m-1)i\\ \alpha\vdash k+(m-1)i\\ \beta\vdash i}} c^{\alpha}_{\gamma,(1^j)}\cdot a^{\gamma/(k-j-i)}_{\beta',(m)}\cdot a^{\hat{\mu}/\alpha}_{\lambda/\beta,(m-1)}\\
		&= \sum_{i=0}^{k-1} (-1)^{k+i} \sum_{j=1}^{k-i} \sum_{\beta\vdash i} (-1)^{j+1}\cdot Y^\beta_{i,j},
	\end{align*}
	where we set
	\begin{align*}
		Y^\beta_{i,j}\coloneqq \sum_{\substack{\gamma\vdash k-j + (m-1)i\\ \alpha\vdash k+(m-1)i}} c^{\alpha}_{\gamma,(1^j)}\cdot a^{\gamma/(k-j-i)}_{\beta',(m)}\cdot a^{\hat{\mu}/\alpha}_{\lambda/\beta,(m-1)}.
	\end{align*}
	Next, we will simplify $Y^\beta_{i,j}$, for any fixed $\beta$, $i$ and $j$. To ease notation, in the rest of this proof we will abbreviate sums over all partitions of a given size. That is, we shorten $\sum_{\omega\vdash t}$ to $\sum_\omega$ (the size $t$ will always be clear from context). We use \eqref{eqn:skew-plethysm} to obtain
	\begin{align*}
		Y^\beta_{i,j}=\sum_{\gamma} \sum_\alpha c^\alpha_{\gamma,(1^j)}\cdot a^{\hat{\mu}/\alpha}_{\lambda/\beta,(m-1)} \cdot \sum_\varepsilon c^\gamma_{\varepsilon,(k-j-i)}\cdot a^\varepsilon_{\beta',(m)}.
	\end{align*}
	By \cite[(2.1)]{SLthesis}, we have that $\sum_\gamma c^\alpha_{\gamma,(1^j)}\cdot c^{\gamma}_{\varepsilon,(k-j-i)} = c^{\alpha}_{\varepsilon,(k-j-i),(1^j)} = \langle \chi^{\alpha/\varepsilon}, \chi^{k-j-i}\boxtimes \chi^{(1^j)}\rangle$. By Theorem~\ref{thm:LR}, $\chi^{k-j-i}\boxtimes \chi^{(1^j)} = \chi^{H(j)}+\chi^{H(j-1)}$ where $H(j)\coloneqq (k-i-j,1^j)$, except if $j=k-i$ then $\chi^\emptyset\boxtimes \chi^{(1^{k-i})}=\chi^{H(k-i-1)}$ only, i.e.~we treat the $\chi^{H(k-i)}$ term as the zero character. Hence
	\begin{align*}
		Y^\beta_{i,j} &= \sum_\alpha \sum_\varepsilon \langle \chi^{\alpha/\varepsilon}, \chi^{H(j)}+\chi^{H(j-1)}\rangle \cdot a^\varepsilon_{\beta',(m)}\cdot a^{\hat{\mu}/\alpha}_{\lambda/\beta,(m-1)}\\
		&= \sum_\alpha \sum_\varepsilon (c^\alpha_{\varepsilon,H(j)} + c^\alpha_{\varepsilon,H(j-1)}) \cdot a^\varepsilon_{\beta',(m)}\cdot a^{\hat{\mu}/\alpha}_{\lambda/\beta,(m-1)}
	\end{align*}
	(where we omit the $c^\alpha_{\varepsilon,H(k-i)}$ term). Since $H(0)=(k-i)$, we obtain
	\[ \sum_{j=1}^{k-i}(-1)^{j+1} \cdot Y^\beta_{i,j} = \sum_\alpha \sum_\varepsilon c^\alpha_{\varepsilon,(k-i)}\cdot a^\varepsilon_{\beta',(m)}\cdot a^{\hat{\mu}/\alpha}_{\lambda/\beta,(m-1)} = \sum_\alpha a^{\alpha/(k-i)}_{\beta',(m)}\cdot a^{\hat{\mu}/\alpha}_{\lambda/\beta,(m-1)}.\]
	Thus, we finally obtain
	\begin{align*} X &= \sum_{i=0}^{k-1}(-1)^{k+i} \cdot \sum_{\beta\vdash i} \sum_{j=1}^{k-i}(-1)^{j+1}\cdot Y^\beta_{i,j} \\&= \sum_{i=0}^{k-1}(-1)^{k+i}\cdot \sum_{\beta\vdash i}\ \sum_{\alpha\vdash k+(m-1)i} a^{\alpha/(k-i)}_{\beta',(m)}\cdot a^{\hat{\mu}/\alpha}_{\lambda/\beta,(m-1)}, \end{align*}
	as desired.
\end{proof}
